\documentclass{article}
\usepackage{graphicx}
\begin{document}
\section{Mapping with an MRRG}

A modulo routing resource graph replicates each resource over time and
connects those copies wherever a value may advance. The graph builder
must encode this distinction, and the mapper must reserve the full
implementation rather than only its result vertex.

However, the MRRG is appealing because time expansion converts temporal
availability into spatial graph topology. For a candidate initiation
interval, it replicates each resource over time.

\begin{figure}[h]
\centering
\includegraphics[width=0.4\linewidth]{fig}
\caption{One ALU's ADD, MUL and NOP behaviours as an illustrative MRRG.}
\end{figure}

Moreover, an MRRG decomposes an undirected hypergraph netlist of
connected components into a directed graph. Multiple directed edges at
the same modulo time may represent connections through the same
physical connection.

\section{The MRRG dilemma}

Even if every routing and architecture-model limitation above
disappeared, the candidate selection process would remain. Choosing the
first candidate at which to construct the MRRG requires architecture
semantics that exist independently of that graph.

These design choices in MRRG-based mappers over the years made it
infeasible to universalize the mapping independently from what the
architectural features are.
\end{document}
